\documentclass[a4paper,12pt]{article}
\usepackage[utf8]{inputenc}
\usepackage[T1]{fontenc}
\usepackage[ngerman]{babel}
\usepackage{amsmath, amssymb}
\usepackage{geometry}
\usepackage{tikz}
\usetikzlibrary{shapes.geometric, arrows.meta, positioning, fit, backgrounds}
\usepackage{parskip}
\usepackage{titlesec}
\usepackage{booktabs}
\usepackage{enumitem}

\geometry{left=3cm, right=3cm, top=2.5cm, bottom=2.5cm}

\titleformat{\section}{\large\bfseries}{}{0em}{}[\titlerule]
\titleformat{\subsection}{\normalsize\bfseries}{}{0em}{}

\tikzset{
  box/.style={rectangle, draw, rounded corners=3pt, text width=4.2cm,
              align=center, minimum height=0.9cm, font=\small, inner sep=4pt},
  boxw/.style={box, text width=5.2cm},
  boxn/.style={box, text width=3.8cm},
  arrow/.style={-{Latex[length=2mm]}, thick},
}

\begin{document}

\begin{center}
  {\LARGE\bfseries Studierhinweise zu Kurseinheit 6}\\[0.5em]
  {\large Analysis (Modul 61211) -- Kurven}
\end{center}

\vspace{1em}

In dieser Kurseinheit stehen zunächst der Kurvenbegriff und die Berechnung
der Länge von Kurven im Mittelpunkt. Dieser Kurvenbegriff ist nicht einfach,
und Sie sollten sich das Vorgehen sehr gut verdeutlichen, da es in mancher
Hinsicht typisch ist: Äquivalente Objekte werden zu einem neuen Objekt
zusammengefasst. Zum Glück steht der hohen Abstraktionsstufe des Kurvenbegriffs
die angenehme Möglichkeit gegenüber, sich Kurven im Raum sehr schön
veranschaulichen zu können.

In Abschnitt~6.3 werden, aufbauend auf dem Kurvenbegriff, Integrale längs
Kurven definiert, einmal für skalare Funktionen, zum anderen (und das ist
der wichtigere Fall) für Funktionen mit Werten in $\mathbb{R}^n$.
Kurvenintegrale spielen in den Anwendungsgebieten (z.\,B.\ in der Physik)
eine große Rolle. Auch innermathematisch sind sie bedeutungsvoll, da sie
direkt mit der Konstruktion von Stammfunktionen (physikalisch: Potenzialen)
zusammenhängen. Die beiden wichtigsten Sätze sind 6.4.4 und 6.4.6; sie
verallgemeinern den Hauptsatz der Differenzial- und Integralrechnung 5.1.10.

Der letzte Abschnitt~6.5 ist geometrischen Anwendungen der Integralrechnung
gewidmet, nämlich der Inhaltsberechnung von einfachen Flächen (im
$\mathbb{R}^2$) und Körpern (im $\mathbb{R}^3$). Hauptsächlich dient der
Abschnitt dazu, noch einmal Integrationstechniken zu üben.

% -----------------------------------------------------------------------
\section*{Struktur der Kurseinheit 6}

\begin{center}
\begin{tikzpicture}[node distance=0.5cm and 0.6cm]

  % Column headers
  \node[font=\bfseries\small] (hA) at (0,0)    {Vorkenntnisse};
  \node[font=\bfseries\small] (hB) at (8.5,0)  {Neue Inhalte};

  % ---- Left column: Vorkenntnisse ----
  \node[box, below=0.8cm of hA] (v1) {2.3--2.5\\Stetigkeit};
  \node[box, below=1.0cm of v1] (v2) {3.3, 3.5--3.6\\Differenzierbarkeit};
  \node[box, below=1.0cm of v2] (v3) {5.1--5.3\\Integration};

  % ---- Right column: Neue Inhalte ----
  \node[box, below=0.8cm of hB] (n1) {6.1 parametrisierte Kurve,\\Äquivalenz};
  \node[box, below=0.3cm of n1] (n2) {6.1 Kurve,\\Parameterdarstellung};
  \node[box, below=0.3cm of n2] (n3) {6.1 entgegengesetzte Kurve,\\Aneinanderfügen};
  \node[box, below=0.5cm of n3] (n4) {6.2 Länge einer Kurve,\\rektifizierbare Kurve};
  \node[box, below=0.3cm of n4] (n5) {6.2 Länge einer stetig\\differenzierbaren Kurve};
  \node[box, below=0.3cm of n5] (n6) {6.2 Bogenlänge,\\ausgez.\ Parameterdarstellung};
  \node[box, below=0.3cm of n6] (n7) {6.2 glatte und\\stückweise glatte Kurven};
  \node[box, below=0.5cm of n7] (n8) {6.3 Kurvenintegral für\\skalare Funktionen und\\Vektorfunktionen};
  \node[box, below=0.5cm of n8] (n9) {6.4 Stammfunktion,\\Integrabilitätsbedingungen};
  \node[box, below=0.3cm of n9] (n10){6.4 Wegunabhängigkeit,\\Existenz von Stammfunktionen};
  \node[box, below=0.5cm of n10](n11){6.5 Flächeninhalte,\\Volumina};

  % Arrows right column (lesson flow)
  \draw[arrow] (n1) -- (n2);
  \draw[arrow] (n2) -- (n3);
  \draw[arrow] (n3) -- (n4);
  \draw[arrow] (n4) -- (n5);
  \draw[arrow] (n5) -- (n6);
  \draw[arrow] (n6) -- (n7);
  \draw[arrow] (n7) -- (n8);
  \draw[arrow] (n8) -- (n9);
  \draw[arrow] (n9) -- (n10);
  \draw[arrow] (n10) -- (n11);

  % Cross arrows (Vorkenntnisse -> neue Inhalte)
  \draw[arrow, dashed] (v1.east) -- ++(0.3,0) |- (n1.west);
  \draw[arrow, dashed] (v2.east) -- ++(0.3,0) |- (n5.west);
  \draw[arrow, dashed] (v3.east) -- ++(0.3,0) |- (n8.west);
  \draw[arrow, dashed] (v3.east) -- ++(0.6,0) |- (n11.west);

\end{tikzpicture}
\end{center}

\newpage

% -----------------------------------------------------------------------
\section*{Zielelement 6.1 -- Der Kurvenbegriff}

\subsection*{Lerninhalte}

\begin{center}
\begin{tikzpicture}[node distance=0.5cm and 1.0cm]
  \node[box] (para) {parametrisierte Kurve\\$\varphi: I \to \mathbb{R}^n$, Spur $\varphi(I)$};
  \node[box, right=1.2cm of para] (aeq) {Äquivalenz parametrisierter Kurven\\$\varphi \sim \psi$, Äquivalenzklasse $[\varphi]$};
  \node[box, below=0.7cm of $(para)!0.5!(aeq)$] (kurve) {Kurve $W$, Parameterdarstellung,\\Spur $|W|$, Anfangs-/Endpunkt};
  \node[box, below=0.5cm of kurve] (ent) {entgegengesetzte Kurve $-W$};
  \node[box, below=0.5cm of ent] (ane) {Aneinanderfügen von Kurven\\$W_1 + W_2$};
  \draw[arrow] (para) -- (aeq);
  \draw[arrow] (para.south) -- (kurve.north west);
  \draw[arrow] (aeq.south) -- (kurve.north east);
  \draw[arrow] (kurve) -- (ent);
  \draw[arrow] (ent) -- (ane);
\end{tikzpicture}
\end{center}

\subsection*{Lernziele}

Nach Durcharbeiten dieses Abschnitts sollten Sie in der Lage sein:
\begin{itemize}
  \item den Kurvenbegriff aus dem Begriff der parametrisierten Kurve und dem
        Äquivalenzbegriff für parametrisierte Kurven zu entwickeln,
  \item für Kurven $W$ und $\widetilde W$ die entgegengesetzte Kurve $-W$
        sowie die zusammengesetzte Kurve $W + \widetilde W$ zu erklären.
\end{itemize}

\subsection*{Selbstkontrollelement 6.1}

Die beiden Kurven $W := [\varphi]$ und $\widetilde W := [\widetilde \varphi]$ mit
\[
  \varphi: [0, 2\pi] \to \mathbb{R}^2,\quad t \mapsto \varphi(t) := \begin{pmatrix}\cos t \\ \sin t\end{pmatrix}
\]
bzw.
\[
  \widetilde\varphi: [0, 2\pi] \to \mathbb{R}^2,\quad t \mapsto \widetilde\varphi(t) := \begin{pmatrix}\cos(2t) \\ \sin(2t)\end{pmatrix}
\]
haben die gleiche Spur, gleichen Anfangs- und Endpunkt sowie gleichen
Durchlaufungssinn. Gilt also $W = \widetilde W$?

\newpage

% -----------------------------------------------------------------------
\section*{Zielelement 6.2 -- Länge einer Kurve}

\subsection*{Lerninhalte}

\begin{center}
\begin{tikzpicture}[node distance=0.5cm and 1.0cm]
  \node[box] (str) {gerichtete Strecke $S(a,b)$,\\einbeschriebener Polygonzug};
  \node[box, below=0.5cm of str] (laenge) {Länge einer Kurve $L(W)$,\\rektifizierbare Kurve};
  \node[box, below=0.5cm of laenge] (difb) {Länge einer stetig\\differenzierbaren Kurve\\(Satz 6.2.5)};
  \node[box, below=0.5cm of difb] (bogen) {Bogenlänge $s_\varphi$,\\ausgezeichnete\\Parameterdarstellung $\Phi$};
  \node[box, below=0.5cm of bogen] (glatt) {glatte und\\stückweise glatte Kurven};
  \draw[arrow] (str) -- (laenge);
  \draw[arrow] (laenge) -- (difb);
  \draw[arrow] (difb) -- (bogen);
  \draw[arrow] (bogen) -- (glatt);
\end{tikzpicture}
\end{center}

\subsection*{Lernziele}

Nach Durcharbeiten dieses Abschnitts sollten Sie in der Lage sein:
\begin{itemize}
  \item die Länge einer Kurve mithilfe des Begriffs des einbeschriebenen
        Polygonzuges zu definieren,
  \item die Länge einer stetig differenzierbaren Kurve mithilfe eines
        Integrals zu berechnen,
  \item die Bogenlänge (bezüglich einer Parameterdarstellung) einer glatten
        Kurve sowie die ausgezeichnete Parameterdarstellung zu definieren und
        an Beispielen zu erläutern.
\end{itemize}

\subsection*{Selbstkontrollelement 6.2}

Der Graph einer stetigen Funktion $f: [\alpha, \beta] \to \mathbb{R}$ mit
$\alpha < \beta$ kann als Spur der durch
\[
  \varphi: [\alpha, \beta] \to \mathbb{R}^2,\quad t \mapsto \varphi(t) := \begin{pmatrix}t \\ f(t)\end{pmatrix}
\]
erzeugten Kurve aufgefasst werden. Können Sie die Länge dieser Kurve durch
ein Integral darstellen, wenn $f$ stetig differenzierbar ist?

\newpage

% -----------------------------------------------------------------------
\section*{Zielelement 6.3 -- Kurvenintegrale}

\subsection*{Lerninhalte}

\begin{center}
\begin{tikzpicture}[node distance=0.5cm and 1.2cm]
  \node[box] (skal) {Kurvenintegral\\für skalare Funktionen\\$\int_W f\,ds$};
  \node[box, right=1.5cm of skal] (vekt) {Kurvenintegral\\für Vektorfunktionen\\$\int_W f \cdot dx$};
  \node[box, below=0.5cm of skal] (bspS) {Beispiele};
  \node[box, below=0.5cm of vekt] (bspV) {Beispiele, Rechenregeln\\(Linearität, Additivität,\\Abschätzung)};
  \draw[arrow] (skal) -- (bspS);
  \draw[arrow] (vekt) -- (bspV);
\end{tikzpicture}
\end{center}

\subsection*{Lernziele}

Nach Durcharbeiten dieses Abschnitts sollten Sie in der Lage sein:
\begin{itemize}
  \item Kurvenintegrale für skalare und für vektorwertige Funktionen zu
        definieren und an Beispielen zu erläutern.
\end{itemize}

\subsection*{Selbstkontrollelement 6.3}

Können Sie
\[
  \int_R (x\,dx + y\,dy)
\]
für die ,,Rechteckskurve'' $R$ (das ist der Rand des Rechtecks
\[
  \left\{\begin{pmatrix}x\\y\end{pmatrix} \in \mathbb{R}^2 \;\middle|\; |x| = 2,\; |y| \leq 1\right\}
  \cup
  \left\{\begin{pmatrix}x\\y\end{pmatrix} \in \mathbb{R}^2 \;\middle|\; |y| = 1,\; |x| \leq 2\right\},
\]
einmal im Gegenuhrzeigersinn durchlaufen, beginnend und endend in
$a := \begin{pmatrix}2 \\ 1\end{pmatrix}$, berechnen?

\newpage

% -----------------------------------------------------------------------
\section*{Zielelement 6.4 -- Stammfunktionen}

\subsection*{Lerninhalte}

\begin{center}
\begin{tikzpicture}[node distance=0.5cm and 1.0cm]
  \node[box] (stamm) {Stammfunktion einer\\vektorwertigen Funktion\\${}^t F'(x) = f(x)$};
  \node[box, right=1.5cm of stamm] (integ) {Integrabilitäts-\\bedingungen\\$D_i f_j = D_j f_i$};
  \node[box, below=0.7cm of stamm] (wegun) {Kurvenintegral von\\Funktionen mit\\Stammfunktion\\$\int_W f \cdot dx = F(b) - F(a)$};
  \node[box, below=0.7cm of integ] (exist) {Existenz und Berechnung\\von Stammfunktionen\\(sternförmiges Gebiet)};
  \draw[arrow] (stamm) -- (integ);
  \draw[arrow] (stamm) -- (wegun);
  \draw[arrow] (integ) -- (exist);
  \draw[arrow] (wegun) -- (exist);
\end{tikzpicture}
\end{center}

\subsection*{Lernziele}

Nach Durcharbeiten dieses Abschnitts sollten Sie in der Lage sein:
\begin{itemize}
  \item den Begriff der Stammfunktion einer vektorwertigen Funktion zu
        formulieren und an Beispielen zu erläutern,
  \item die Bedeutung einer Stammfunktion für die Berechnung von
        Kurvenintegralen zu erläutern,
  \item hinreichende Bedingungen für die Existenz von Stammfunktionen sowie
        ein Berechnungsverfahren anzugeben.
\end{itemize}

\subsection*{Selbstkontrollelement 6.4}

Können Sie mit den Ergebnissen dieses Abschnitts erläutern, warum das
Integral im Selbstkontrollelement~6.3 den Wert $0$ haben muss?

\newpage

% -----------------------------------------------------------------------
\section*{Zielelement 6.5 -- Flächen- und Volumenberechnungen}

\subsection*{Lerninhalte}

\begin{center}
\begin{tikzpicture}[node distance=0.5cm and 1.0cm]
  \node[boxw] (flae) {Berechnung des Inhalts einer Fläche\\zwischen zwei Funktionsgraphen};
  \node[boxw, below=0.5cm of flae] (koerp) {Berechnung des Volumens eines Körpers\\zwischen zwei Funktionsgraphen};
  \node[boxw, below=0.5cm of koerp] (rot) {Berechnung des Volumens\\eines Rotationskörpers};
  \draw[arrow] (flae) -- (koerp);
  \draw[arrow] (koerp) -- (rot);
\end{tikzpicture}
\end{center}

\subsection*{Lernziele}

Nach Durcharbeiten dieses Abschnitts sollten Sie erklären können:
\begin{itemize}
  \item wie der Flächeninhalt zwischen den Graphen zweier stetiger Funktionen
        zu berechnen ist,
  \item wie das Volumen eines dreidimensionalen Körpers zwischen den Graphen
        zweier stetiger Funktionen zu berechnen ist, insbesondere das eines
        Rotationskörpers.
\end{itemize}

\subsection*{Selbstkontrollelement 6.5}

Sei $Q$ ein Rechteck
$\left\{\begin{pmatrix}x\\y\end{pmatrix} \in \mathbb{R}^2 \;\middle|\; a \leq x \leq b,\; c \leq y \leq d\right\}$
mit $a < b$, $c \leq d$, und sei $f: Q \to \mathbb{R}$ eine stetige Funktion
mit $f(x,y) \geq 0$ für $\begin{pmatrix}x\\y\end{pmatrix} \in Q$. Können Sie
das ,,iterierte Integral''
\[
  \int_a^b \left( \int_c^d f(x,y)\,dy \right) dx
\]
als Volumen eines dreidimensionalen Körpers interpretieren?

\end{document}
